\documentclass[letterpaper]{article}
\usepackage[submission]{aaai2027}
\usepackage[hyphens]{url}
\usepackage{graphicx}
\urlstyle{rm}
\def\UrlFont{\rm}
\usepackage{natbib}
\usepackage{caption}
\frenchspacing
\usepackage{amsmath}
\usepackage{amssymb}
\usepackage{amsthm}
\usepackage{algorithm}
\usepackage{algpseudocode}   % do NOT also load `algorithmic' (clash)
\usepackage{booktabs}

\graphicspath{{Figures/}}

% AAAI sections are unnumbered (secnumdepth=0), so theorem environments
% must NOT use the [section] counter, and sections are referred to by
% name, never by \ref.
\newtheorem{proposition}{Proposition}
\newtheorem{remark}{Remark}

\pdfinfo{
/TemplateVersion (2027.1)
}

\setcounter{secnumdepth}{0}

\title{TempoMem: Persistent Scene Memory for Efficient Video Transformer Inference}

% Submission is anonymous; author block added on acceptance.

\begin{document}

\maketitle

\begin{abstract}
Transformer-based models have demonstrated outstanding performance in video understanding tasks due to their capacity to capture long-range dependencies. However, their high computational cost, combined with the massive volume of streaming video data, poses significant challenges for real-time deployment on resource-constrained edge devices. Existing efficiency methods reduce redundancy by comparing consecutive frames, but discard the wealth of information already observed earlier in the stream: once a token falls outside the immediate temporal window, it must be recomputed from scratch when the same scene content reappears. To address this limitation, we propose TempoMem, a training-free framework that exploits long-term recurrence in streaming video. Through a novel clustering algorithm, TempoMem incrementally builds a Persistent Scene Memory (PSM), a compact key-value representation of previously observed scene content, and reuses it to bypass computation for recurring tokens. Because it is training-free and architecture-agnostic, TempoMem is designed to be layered on top of existing state-of-the-art acceleration methods, compounding their savings rather than competing with them. When layered on top of an existing state-of-the-art acceleration method, experimental results show that TempoMem reduces computation by 75\% compared to the base model and by a further 49\% in steady state compared to the accelerated method alone, at a cost of only 2.5\% in Top-1 accuracy, improving the accuracy-efficiency tradeoff for video transformers in streaming settings.
\end{abstract}

\section{Introduction}

Vision Transformers (ViTs) \cite{vit} represent images as token sequences
and have become the dominant backbone for video understanding
\cite{vivit,videomae}. Their cost, however, scales with the number of
tokens processed, and video multiplies that number by the frame rate.
This cost collides with a growing deployment reality: cameras on
vehicles, drones, and IoT installations produce continuous streams that
must be analyzed close to the sensor, on hardware whose compute and
energy budgets are orders of magnitude below a datacenter GPU. For such
edge settings the relevant workload is not a short clip but an unbounded
stream, and the relevant question is not how fast one frame can be
processed but how little work each additional frame can be made to cost.

A rich body of work attacks this redundancy. Spatial token-reduction
methods prune, select, or merge less informative tokens within a single
frame \cite{tome,evit,dynamicvit,avit,ats}; temporal methods reuse or
selectively update information between frames, recomputing only tokens
that changed \cite{eventful,stgt,deltacnn}, masking background regions
\cite{maskvd}, or cascading several such gates \cite{cstvit}. Despite
their diversity, these temporal methods share one structural trait:
every frame, each of them partitions the tokens into a set worth
recomputing and a set whose previous result can be reused. The decision
is made per frame, against the immediately preceding frames, and
whatever was learned about the scene earlier in the stream is discarded.
Yet most of what a video shows is not merely similar to the last frame;
it is \emph{persistent}. Roads, walls, floors, and static objects remain
visually stable for hundreds of frames, disappear behind the camera, and
reappear. A per-frame decision rule pays for such content again and
again. We therefore ask a different question: \emph{which visual
computations should not need to be performed again at all?}

We answer it with \textbf{TempoMem}, a training-free, memory-centric
inference framework. Its central object is a \emph{Persistent Scene
Memory} (PSM): a compact, per-block summary of the scene content
observed so far, built by a clustering pass that merges the stable
tokens of a warm-up window into a bounded set of prototypes, stored
together with their exact key/value projections. On later frames,
incoming tokens are matched against the memory by content; a matched
token reuses its prototype's cached projections and can be represented
by the prototype in the sequence itself, so the model computes only what
the memory cannot already explain. Because matching is by content rather
than by grid position, background that shifted under camera motion can
still hit the memory, a signal complementary to the per-position change
detection of existing temporal methods.

TempoMem operates in two modes that share the same memory but serve
different deployments. \emph{Standalone} TempoMem is a complete
accelerator: it splits each frame into a stable background and a
task-critical foreground using an attention-based saliency score,
caches and reuses only the background, and recomputes the foreground
exactly at every frame. \emph{Layered} TempoMem instead attaches to an
existing temporal token-reduction method. Since all such methods already
flag tokens as useful or useless each frame, the layered mode simply
takes the tokens the host method classifies as useful, checks them
against the PSM, and removes those the memory already explains. It
never adds tokens back, so it can only subtract work from its host, and
with the memory disabled it reduces exactly to the host method. This
makes TempoMem a compounding accelerator rather than a competing one:
any present or future gating method becomes a candidate host. Both
modes are attractive at the edge for the same underlying reason: the
memory is bounded regardless of stream length, its construction cost is
paid once and amortized over the stream, and the reuse rate acts as a
dial that trades a controlled amount of accuracy for large reductions
in per-frame computation.

Our contributions are:
\begin{enumerate}
    \item \textbf{TempoMem}, a training-free persistent-memory framework
    for video transformer inference: each block maintains a Persistent
    Scene Memory of prototypes built by clustering stable scene content,
    and couples prototype-based key/value reuse with token substitution.

    \item \textbf{Two operating modes.} A standalone accelerator driven
    by an attention-based foreground/background split, and a layered
    mode that filters the recompute set of any existing temporal
    token-reduction method, strictly shrinking its work and reducing to
    the unmodified host when disabled.

    \item \textbf{A cost model for streaming inference.} The clustering
    schedule bounds the memory independently of stream length, a
    break-even analysis gives the condition under which matching pays
    for itself, and an amortized accounting captures the long-stream
    regime that edge deployments actually run in.

    \item \textbf{An honest empirical scoping of content-based reuse.}
    Layered on a representative gating accelerator, TempoMem discards a
    large majority of the recompute set on classification at a small,
    one-time accuracy cost, while on dense detection the same filter
    trades accuracy much less favorably. We analyze why the prediction
    readout, global versus dense, determines when content-based reuse
    helps, a finding we believe is useful beyond this specific method.
\end{enumerate}

\section{Related Work}

\paragraph{Spatial token reduction.}
A first family accelerates ViTs by shrinking the token sequence within a
single image or frame. DynamicViT \cite{dynamicvit} and A-ViT \cite{avit}
learn to prune uninformative tokens; EViT \cite{evit} keeps the tokens
most attended by the class token; ATS \cite{ats} samples tokens
adaptively by importance; ToMe \cite{tome} merges similar tokens with
bipartite soft matching, requiring no training. These methods reduce
intra-frame redundancy and treat every frame as an independent inference
problem. TempoMem borrows ToMe's bipartite merging as the clustering
subroutine of its memory, but applies it \emph{across} a temporal window
to build a reusable cache rather than within one frame. The two act on
disjoint token sets in the standalone mode, where a single-frame reducer
can handle the foreground while TempoMem handles the background, so they
compose rather than compete.

\paragraph{Temporal redundancy.}
A second family reuses computation between frames. DeltaCNN
\cite{deltacnn} propagates sparse frame differences through
convolutions; STGT \cite{stgt} gates transformer computation with
learned temporal similarity; Eventful Transformers \cite{eventful}
maintain per-token reference buffers and recompute only the most-changed
tokens; MaskVD \cite{maskvd} masks background regions using detection
priors; CST-ViT \cite{cstvit} cascades direct, offset, and spatial gates
to jointly remove several forms of redundancy. Seen from TempoMem's
perspective, these methods are not only baselines but potential
\emph{hosts}: each of them emits, every frame, a partition of the tokens
into a recompute set and a reuse set, keyed on how much each position
changed recently. That per-frame, position-keyed decision forgets the
stream's history and breaks under camera motion, where every position
changes while scene content largely does not. Layered TempoMem plugs in
exactly at this interface: it takes the host's recompute set, matches it
by content against a persistent memory of the scene, and removes the
tokens the memory already explains. Our evaluation instantiates this
layering on Eventful Transformers as a representative host.

\paragraph{Memory-augmented video transformers.}
MeMViT \cite{memvit} caches compressed activations of past clips to
extend temporal context, and Token Turing Machines \cite{ttm} learn an
external token memory for sequential decision tasks. These works use
memory to improve \emph{accuracy} via longer context and require
training the memory mechanism. TempoMem shares the memory-centric
perspective but uses memory for \emph{efficiency}: it caches key/value
projections of persistent content so they need not be recomputed, is
entirely training-free, and leaves model weights untouched. The
mechanism is loosely analogous to K/V caching in autoregressive language
models, transplanted to the temporal axis of video: previously computed
keys and values are reused rather than recomputed, except that video
requires deciding \emph{which} tokens are stable enough to reuse.

\section{Methodology}
\label{sec:method}

TempoMem rests on one premise: most of what a video shows has been seen
before, so most of what a ViT computes has been computed before.
The framework makes this reuse explicit. During a short warm-up window,
each transformer block summarizes the temporally stable tokens it
observes into a small set of \emph{prototypes} and stores them, together
with their key/value (K/V) projections, in a per-block cache, the
Persistent Scene Memory of the Introduction; in this section we simply
call it the cache. On every later frame, the block matches incoming
tokens against the cache: a matched token \emph{reuses} the prototype's
cached K/V, skipping its own projections, and can additionally be
\emph{substituted} by the prototype in the sequence itself, so attention
and the MLP run on fewer tokens. TempoMem modifies no weights and
requires no fine-tuning. It runs either standalone, where it decides
itself which tokens are eligible for reuse, or layered on a host
accelerator, where the host's own token selection provides that
decision. Both modes share the machinery of this section; they differ
only in how the \emph{cacheable set}, the tokens eligible for caching
and reuse, is chosen. Table~\ref{tab:symbols} collects the symbols used
throughout.

\begin{table}[t]
\centering
\small
\begin{tabular}{ll}
\toprule
Symbol & Definition \\
\midrule
$N$, $C$ & Tokens per frame; token dimensionality \\
$L$ & Number of transformer blocks \\
$N_c$ & Size of the cacheable set in a frame \\
$W$ & Warm-up window length (frames) \\
$T$, $\rho$ & Merge passes; per-pass merge fraction \\
$M$ & Number of prototypes in the cache \\
$E$, $\hat E$ & Pre-LN and post-LN prototype views \\
$\tilde K$, $\tilde V$ & Cached (merged) key/value projections \\
$r$ & Match ratio (standalone reuse budget) \\
$\gamma$ & Reuse fraction of the layered filter \\
$\beta$ & Background ratio of the saliency split \\
$P$ & Refresh period (frames) \\
$M_u$ & Distinct prototypes engaged in a frame \\
$N'$ & Reduced sequence length after substitution \\
$\mathrm{LN}$ & Input LayerNorm of a block \\
$W_Q, W_K, W_V$ & Query, key, value projections \\
\bottomrule
\end{tabular}
\caption{Summary of symbols used in this paper.}
\label{tab:symbols}
\end{table}

A ViT with $L$ blocks processes the stream frame by frame; each block
receives $N$ tokens of dimension $C$ (we omit the batch axis). TempoMem
has four knobs: the warm-up length $W$, the merge schedule $(T, \rho)$,
the reuse budget ($r$ standalone, $\gamma$ layered), and the refresh
period $P$. All cache state is per block.

\subsection{Prototype Cache Construction}
\label{sec:cache}

In streaming video, the cacheable tokens change little across nearby
frames, and many are also mutually similar within a frame. Storing them
verbatim over a $W$-frame window would grow the cache linearly with $W$,
so TempoMem compresses everything accumulated during the window into a
bounded set of prototypes and stores each prototype together with its
identically merged K/V. The compression primitive is bipartite soft
matching (BSM) \cite{tome}: one pass splits $n$ tokens into two
alternating sets, links each token in the first set to its most
cosine-similar token in the second, keeps the strongest links, and
averages each linked pair into a single token. A single pass compresses
by at most a factor of two, which is not enough to summarize a
multi-frame token set, so we compose $T$ passes, each removing a
fraction $\rho$ of the surviving tokens. The survivors, running averages
of everything merged into them, are our prototypes. We make no
optimality claim for this greedy clustering; its job is to bound the
cache size while keeping every prototype close, in cosine distance, to
the tokens it absorbed, which is exactly the property the matching step
relies on.

\begin{proposition}[Cache size bound]
\label{prop:cachesize}
If $n_0$ tokens are accumulated over the warm-up window, then after $T$
passes with merge fraction $\rho$ the cache holds
$M = \lceil n_0 (1-\rho)^T \rceil$ prototypes, up to per-pass rounding.
Doubling the window doubles $n_0$ but is offset by a single extra pass
whenever $\rho = 0.5$, so the final cache size is effectively decoupled
from the window length.
\end{proposition}

During warm-up the block runs unmodified and appends its cacheable
tokens and their exact K/V to a buffer. When the window closes, the
merge is computed once on the pooled buffer, using post-LN features as
the similarity metric, and the same merge is applied to three aligned
tensors: pre-LN tokens, post-LN tokens, and K/V. The cache therefore
holds four $M \times C$ matrices: a pre-LN prototype view $E$, a post-LN
view $\hat E$, and the merged projections $\tilde K, \tilde V$. Merging
once over the pooled window, rather than per frame, is what makes
Proposition~\ref{prop:cachesize} apply to the final cache rather than to
each frame's contribution. Two views are kept because matching and the
query projection consume post-LN features while the residual stream
consumes pre-LN features; caching a single view and reusing it for both
purposes compares vectors from different distributions during matching
and feeds $W_Q$ activations it never sees in the base model. One
approximation should be kept in mind: because the merge is also applied
to K/V, matched queries later attend to \emph{averaged} keys and values.
Attention is nonlinear in the keys, so this approximates attending to
the original key set; the error is governed by how tight each
prototype's group is, which the similarity-ordered BSM links keep small.

\subsection{Cache-Aware Recomputation}
\label{sec:recompute}

At matching frames, the block uses the cache to avoid work
(Algorithm~\ref{alg:standalone} instantiates the full procedure for the
standalone mode and Algorithm~\ref{alg:layered} for the layered mode).
The frame's tokens split into the cacheable set, containing $N_c$
tokens, and the remaining tokens, which the pipeline always recomputes
exactly. Each cacheable token is compared with every prototype by cosine
similarity on post-LN features and assigned its single best-matching
prototype; the $\lfloor r N_c \rfloor$ tokens with the highest
best-match similarity are declared \emph{matched}, and the rest stay
unmatched. Since many tokens often map to the same prototype, duplicate
assignments are collapsed, leaving
$M_u \le \min(\lfloor r N_c \rfloor, M)$ distinct prototypes in play.
Using a fixed count rather than a similarity threshold keeps tensor
shapes batch-uniform and makes the compute budget directly tunable; the
price is that easy and hard frames are forced to the same reuse rate
(see the Discussion and Limitations section).

Matched tokens never enter the key and value projections: their
contribution to attention is represented by the cached
$\tilde K, \tilde V$ rows of their prototypes, and fresh K/V are
computed only for the unmatched cacheable tokens and the
always-recomputed tokens. Matched tokens can additionally be removed
from the sequence itself and represented by their prototypes. Writing
$E_{\mathrm{m}}$ for the pre-LN views of the $M_u$ distinct matched
prototypes, $X_{\mathrm{u}}$ for the unmatched cacheable tokens, and
$X_{\mathrm{a}}$ for the always-recomputed tokens, the block proceeds on
the reduced sequence
\begin{equation}
  X' \;=\; \big[\, E_{\mathrm{m}} \,\big\|\, X_{\mathrm{u}} \,\big\|\,
  X_{\mathrm{a}} \,\big],
  \qquad N' = |X'| \;\le\; N .
  \label{eq:reduced}
\end{equation}
The pre-LN prototype view enters the residual stream and the post-LN
view enters $W_Q$. Queries are computed for all $N'$ tokens; attention,
the output projection, and the MLP then run on the reduced sequence, so
savings are quadratic in the attention and linear in every pointwise
stage. Because deduplication collapses all tokens matching the same
prototype into one sequence element, $N'$ is typically well below $N$;
the per-frame cost is data-dependent, and the counts we report are
measured, not worst-case.

\begin{algorithm}[t]
\caption{Standalone TempoMem (one block, one frame)}
\label{alg:standalone}
\begin{algorithmic}[1]
\Require
  Frame tokens; phase (\textsc{caching} or \textsc{matching});
  cache $(E, \hat E, \tilde K, \tilde V)$; buffer $B$;
  ratios $\beta$, $r$
  \vspace{0.3em}

\If{phase is \textsc{caching}}
  \Comment{Warm-up or refresh frame}
  \State Run the block unmodified (dense attention and MLP)
  \State Score tokens with the saliency of Eq.~\ref{eq:saliency};
    the $\lfloor \beta N \rfloor$ lowest scores form the background
  \State Append the background tokens and their exact K/V to $B$
  \If{this is the last caching frame}
    \State Merge $B$ with $T$ BSM passes into
      $(E, \hat E, \tilde K, \tilde V)$; clear $B$
  \EndIf
  \State \Return the dense block output
\Else
  \Comment{Matching frame}
  \State Split the normalized tokens into background $X_b$ and
    foreground $X_f$ via Eq.~\ref{eq:saliency}
  \State Assign each background token its most cosine-similar prototype
  \State Mark the $\lfloor r N_b \rfloor$ best-matching tokens as matched;
    $E_{\mathrm{m}} \gets$ their distinct prototypes
  \State $X' \gets [\,E_{\mathrm{m}} \mathbin{\Vert} X_{\mathrm{u}}
    \mathbin{\Vert} X_f\,]$
    \hfill\Comment{Token substitution}
  \State $Q \gets$ query projection of all $N'$ tokens of $X'$
  \State $K \gets [\,\tilde K_{\mathrm{m}} \mathbin{\Vert}
    \text{fresh keys of } X_{\mathrm{u}}, X_f\,]$
    \hfill\Comment{KV reuse}
  \State $V \gets [\,\tilde V_{\mathrm{m}} \mathbin{\Vert}
    \text{fresh values of } X_{\mathrm{u}}, X_f\,]$
  \State \Return block output on $(X', Q, K, V)$
    \hfill\Comment{Attention, projection, MLP on $N'$ tokens}
\EndIf
\end{algorithmic}
\end{algorithm}

\subsection{The TempoMem Pipeline}
\label{sec:pipeline}

TempoMem targets the streaming setting: frames arrive sequentially, past
frames cannot be revisited, and per-frame latency matters more than
amortized-batch throughput. The pipeline cycles through three phases.
The first $W$ frames are \emph{warm-up caching} frames: they run the
unmodified model while every block accumulates its cacheable tokens and
their K/V, so warm-up produces exact outputs and loses no accuracy; the
cache is finalized lazily, with one merge over the pooled window, on the
first matching frame. All subsequent frames are \emph{matching} frames
and run the matching branch of Algorithm~\ref{alg:standalone} (or
Algorithm~\ref{alg:layered}) in every TempoMem block. Finally, reused
prototypes age: as the scene evolves, tokens that still match the cache
are represented by increasingly stale features, and after a scene cut
the cache may describe content that no longer exists. Every $P$ frames
the pipeline therefore re-enters caching mode and rebuilds the cache,
bounding staleness by $P$ frames at the cost of one full-model pass
whose overhead amortizes to a factor $1/P$ (see the Complexity Analysis
section).
% [TBD] Strongly recommended: add a one-column pipeline overview figure
% here (warm-up caching -> merge to prototypes -> matching frames with
% KV reuse + substitution -> periodic refresh), referenced as Fig. 1.
% AAAI reviewers expect a method diagram; the paper currently has none.

Unless stated otherwise, we report \emph{steady-state} per-frame cost:
matching frames only, with refresh amortized. Warm-up cost, which occurs
once per stream (or per refresh), is reported separately. This is the
appropriate accounting for long-running streams, and we state it
explicitly to keep comparisons with single-frame methods fair.

\subsection{Standalone TempoMem}
\label{sec:standalone}

In the standalone mode, the cacheable set is the \emph{background}: the
temporally stable regions that dominate most video frames but contribute
little task-relevant information. Foreground tokens, which carry the
motion and object evidence the task depends on, are never matched or
substituted; they are recomputed exactly at every frame. This separation
is what lets TempoMem reuse aggressively without touching the tokens
that matter most. The split is computed from the model's own attention:
background tokens attend diffusely over the frame, while foreground
tokens concentrate their attention on a few relevant tokens. Averaging
the attention map of the previous forward pass over heads to obtain
$\bar A$, token $i$ receives the negative entropy of its attention
distribution,
\begin{equation}
  s_i \;=\; \sum_{j} \bar A_{ij} \log \bar A_{ij},
  \label{eq:saliency}
\end{equation}
normalized to $[0,1]$ across the frame's tokens. The
$\lfloor \beta N \rfloor$ lowest-scoring tokens form the background (the
cacheable set) and the rest the foreground, where $\beta$ is the
background ratio; a fixed split size keeps the partition exact and
batch-uniform. This scoring belongs to the standalone mode only; the
layered mode never computes it, because the host method already
designates the candidate tokens. The attention map is one step stale (it
comes from the previous forward pass through this block); the Discussion
explains why this lag is tolerable for the background side of the
split. Because foreground and background are handled by disjoint
mechanisms, the foreground can additionally be reduced by an
off-the-shelf single-frame reducer such as ToMe \cite{tome} without
double-processing any token.

Token substitution changes sequence length and replaces a token at one
grid position with a prototype averaged over other positions and frames.
This is incompatible with windowed attention using relative positional
embeddings (e.g., ViTDet), where both the window partition and the
relative-position bias require the full, position-faithful grid; after
the first block, positional information is also baked into token
features, so a substituted prototype carries a smeared positional
signature. For such architectures we use a restricted
\emph{KV-reuse-only} variant: the token grid is kept intact, matching is
performed per position against a keyframe cache, and matched positions
skip only their K/V projections while attention runs over the full
window, so partitioning and relative-position biases remain exact. This
variant retains the linear projection savings but not the quadratic
attention savings; we report the two variants separately.

\subsection{Layered TempoMem}
\label{sec:layered}

Temporal token-reduction methods, whatever their internal design,
converge on the same per-frame interface: each frame, the method flags a
subset of tokens as \emph{useful}, to be recomputed, and treats the rest
as \emph{useless}, reusing buffered results for them. Eventful
Transformers \cite{eventful} produce this partition with change-scoring
gates, MaskVD \cite{maskvd} with detection-driven region masks, CST-ViT
\cite{cstvit} with a cascade of temporal and spatial gates. The layered
mode of TempoMem attaches at exactly this interface and asks one further
question: \emph{of the tokens the host considers useful, which ones does
the scene memory already explain?}

Concretely, the host proposes its recompute set of $k$ candidate tokens;
TempoMem matches each candidate against the prototype cache by cosine
similarity and drops the fraction $\gamma \in [0,1)$ of candidates most
similar to a prototype (tokens with a special role, such as a class
token, are always protected). The surviving set is handed back to the
host, which recomputes it through its own machinery; the dropped tokens
fall back on the host's existing reuse path, exactly as if the host had
never flagged them (Algorithm~\ref{alg:layered}). The cache itself is
built from the candidate tokens of the warm-up frames, exactly the
feature distribution the host's selection operates on, and no saliency
computation is involved. Three properties follow directly from this
construction. First, \emph{strict subtraction}: the filter only ever
removes tokens from the host's recompute set, so per-frame FLOPs are
bounded by the host's, up to the matching overhead of order $kMC$, which
is small because both $k$ and $M$ are far below $N$. Second,
\emph{exact reduction}: with $\gamma = 0$ the pipeline is bit-for-bit
the unmodified host, which gives perfectly controlled comparisons.
Third, \emph{host-agnosticism}: nothing in the filter depends on how the
host chose its candidates, so the same implementation layers onto any
method exposing the useful/useless interface.

One refinement matters for models whose prediction is read from a global
token. Dropping a matched candidate means its K/V silently stay stale,
and for class-token readouts the reuse decision never reaches the
output. The \emph{substitution variant} therefore lets the host
recompute the full candidate set, but overwrites each matched
candidate's input with its best-matching prototype, so the token's K/V
are refreshed to the cached summary while the later stages still process
only the kept set. This trades back part of the input-stage saving for
accuracy and is no longer FLOP-dominated by the plain host at that
stage; we quantify the trade in the Evaluation.

\begin{algorithm}[t]
\caption{Layered TempoMem (one block, one frame)}
\label{alg:layered}
\begin{algorithmic}[1]
\Require
  Frame tokens; host method with its per-frame token selection;
  prototype cache $E$; buffer $B$; reuse fraction $\gamma$
  \vspace{0.3em}

\State $C \gets$ recompute set of $k$ tokens proposed by the host
  \hfill\Comment{Tokens the host flags as useful}
\If{phase is \textsc{caching}}
  \State Append the tokens of $C$ to $B$; on the last caching frame,
    merge $B$ with $T$ BSM passes into $E$
  \State \Return standard host output on $C$
\Else
  \State Compare each candidate to every prototype in $E$
    by cosine similarity
  \State $D \gets$ the $\lfloor \gamma k \rfloor$ candidates most similar
    to a prototype
    \hfill\Comment{Special tokens always protected}
  \State $R \gets C \setminus D$
    \hfill\Comment{Kept set: tokens the cache cannot explain}
  \If{substitution variant}
    \State Host recomputes all of $C$, with each input in $D$
      overwritten by its best-matching prototype
  \Else
    \State Host recomputes only $R$
  \EndIf
  \State Dropped tokens follow the host's existing reuse path
  \State \Return host output
\EndIf
\end{algorithmic}
\end{algorithm}

\subsection{Complexity Analysis}
\label{sec:complexity}

Consider one block with MLP expansion $4$; the dense cost per frame is
\begin{equation}
  F_{\text{dense}} \;\approx\; \underbrace{4NC^2}_{\text{QKV+proj}}
  + \underbrace{8NC^2}_{\text{MLP}} + \underbrace{4N^2C}_{\text{attention}} .
\end{equation}
A TempoMem matching frame (the matching branch of
Algorithm~\ref{alg:standalone}) runs every stage on the reduced length
$N'$, computes fresh K/V only for the $N' - M_u$ non-prototype tokens, and
adds the matching product:
\begin{equation}
\begin{aligned}
  F_{\text{TM}} \;\approx\;\;
  & \underbrace{N'C^2}_{Q}
  + \underbrace{2(N'{-}M_u)\,C^2}_{K,V}
  + \underbrace{N'C^2}_{\text{proj}}
  + \underbrace{8N'C^2}_{\text{MLP}} \\
  & + \underbrace{4N'^2C}_{\text{attention}}
  + \underbrace{N_c M C}_{\text{matching}} .
\end{aligned}
\end{equation}
Savings are linear in every pointwise stage (a factor $N'/N$) and quadratic
in attention (a factor $(N'/N)^2$).

\paragraph{Break-even condition.}
The matching overhead pays for itself against the K/V projections alone:
matching costs $N_c M C$, and skipping the matched tokens' projections
saves $2 r N_c C^2$, so matching is profitable whenever
\begin{equation}
  M \;<\; 2\,r\,C .
  \label{eq:breakeven}
\end{equation}
With $C = 768$ and $r = 0.75$, any cache below $1152$ prototypes recoups
its matching cost from the projections alone, and every attention,
projection, and MLP saving is then a net gain. In practice, $M$ sits one to
two orders of magnitude below this bound by
Proposition~\ref{prop:cachesize}.

\paragraph{Amortized stream cost.}
Over a refresh period, the per-frame cost is
$\tfrac{1}{P} F_{\text{dense}} + \tfrac{P-1}{P} F_{\text{TM}}$, plus the
one-time warm-up $W \cdot F_{\text{dense}}$ per stream. The longer the
stream, the smaller the share of the fixed costs: this is the regime a
continuously running camera operates in, and it is where the memory's
construction cost disappears into the noise.

\paragraph{Memory.}
The cache stores $E, \hat E, \tilde K, \tilde V$, i.e., $4MC$ floats per
block ($O(LMC)$ in total), independent of stream length. This replaces the
$O(W N C)$ footprint of verbatim window caching.

\section{Evaluation}
\label{sec:eval}

We evaluate TempoMem on two video transformer regimes: action recognition
with ViViT-B on Kinetics-400, where the prediction is read from a class
token, and video object detection with ViTDet-B on ImageNet VID, where the
prediction depends on the full spatial token grid. The two regimes stress
the method in opposite ways, and we report both findings as measured:
content-based reuse is highly effective for class-token readouts, while
on dense spatial prediction it trades accuracy for compute much less
favorably, with mAP@50 falling as the reuse rate grows. The section
``When Does Content-Based Reuse Help?'' analyzes why.

\subsection{Experimental Setup}
\label{sec:setup}

\paragraph{Models and datasets.}
For action recognition we use ViViT-B \cite{vivit} (factorized
spatio-temporal attention, 16-frame clips at $224 \times 224$, tubelet
$2 \times 16 \times 16$, class-token readout) on Kinetics-400 \cite{k400}.
For detection we use ViTDet-B \cite{vitdet} on ImageNet VID at
$672 \times 672$; its windowed attention with relative-position biases
means TempoMem runs in the KV-reuse-only variant on windowed blocks, as
described in the Methodology. All weights are frozen, pretrained
checkpoints; TempoMem and every compared method are applied as
training-free, inference-time modifications, which isolates the effect of
the reuse mechanism itself.

%% PRELIM: this paragraph states the current evaluation scale explicitly.
%% Replace subset numbers with full-validation results when available.
\paragraph{Evaluation scale.}
Evaluation subsets (100 Kinetics-400 validation clips; 25 ImageNet VID
videos) are sampled once and shared by every method and configuration, so
all within-table comparisons are paired. \textbf{[TBD: full-validation
runs (19{,}877 Kinetics-400 clips; all 639 VID videos) are in progress
and will replace the subset numbers. Completed so far: the layered
$\gamma = 0$ and $\gamma = 0.5$ configurations on full Kinetics-400
(62.4\% and 59.9\% Top-1 at 27.2 and 13.9 matching-frame GFLOPs per
frame, confirming their subset estimates), Eventful at token budget 48
on full Kinetics-400 (67.5\% Top-1 at 53.8 GFLOPs per frame), and one
layered configuration on full VID, reported in the ImageNet VID
section.]}

\paragraph{Baselines.}
\emph{Dense} is the unmodified backbone. \emph{Eventful}
\cite{eventful} recomputes the top-$k$ most-changed tokens per frame and
reuses buffered results for the rest; $k$ is its compute budget.
\emph{STGT} \cite{stgt} is a spatio-temporal token-gating baseline.
\emph{MaskVD} \cite{maskvd} skips background tokens using detection
masks.

\paragraph{Metrics.}
Kinetics-400 reports Top-1 and Top-5 accuracy; ImageNet VID reports
mAP@50. Compute is reported as GFLOPs, measured with in-code
operation counters (instrumented linear and matmul modules) rather than
analytic estimates, so all reductions reflect the operations actually
executed, including matching overhead. Latency is CUDA-event wall-clock
time per frame and memory is peak GPU allocation, both on a single NVIDIA
Quadro RTX 6000 (24\,GB) under a shared warm-up and reset protocol.

\paragraph{Protocol and cost accounting.}
Each video runs the TempoMem pipeline: the first $W$ frames are caching
frames ($W = 4$ for ViTDet, the leading frames of each clip for ViViT),
and all later frames are matching frames. Unless a table says otherwise,
accuracy and GFLOPs are computed \emph{on matching frames only}, for every
method, with baselines evaluated on the same frames; this is the
steady-state accounting defined in the Methodology, and warm-up cost is
reported separately where relevant. Whole-stream numbers (all frames,
warm-up included) are labeled explicitly, and \textbf{we never compare a
matching-only number against a whole-stream number in the same table}.

\subsection{Instantiating the Layered Mode on Eventful Transformers}
\label{sec:eval-instantiation}

The Methodology defines the layered mode against a generic host; here we
make it concrete on Eventful Transformers, the representative host of
this evaluation. An Eventful block gates the qkv input, the
attention-output projection, and the MLP: each gate maintains a
reference copy of its input, scores each token by how much it has
changed since that reference, selects the top-$k$ most-changed tokens
for recomputation, and scatters the recomputed results into an output
accumulator that retains stale values elsewhere. The layered filter
plugs in at the qkv gate: the gate's top-$k$ proposal is the host
recompute set of Algorithm~\ref{alg:layered}, the filter drops the
fraction $\gamma$ of it that best matches the cache (the class token is
always protected), and the surviving set is forced onto all three gates,
whose accumulators supply the reuse path for the dropped tokens. The
cache is built from the gated tokens of the warm-up frames. With
$\gamma = 0$ every gate receives exactly its Eventful proposal, so the
pipeline is Eventful itself, and every comparison below is controlled
down to the single knob $\gamma$.

\subsection{Layered TempoMem on Kinetics-400}
\label{sec:eval-eventful}

Table~\ref{tab:eventful-vivit} sweeps the reuse fraction $\gamma$ on
ViViT-B: every row differs from the first row only in $\gamma$.

\begin{table}[t]
\centering
\small
\begin{tabular}{lccc}
\toprule
$\gamma$ & Top-1 (\%) & Top-5 (\%) & GFLOPs/frame \\
\midrule
0.00 (= Eventful) & 61.0 & 82.0 & 27.2 \\
0.25 & 59.0 & 79.0 & 20.6 \\
0.50 & 59.0 & 79.0 & 13.9 \\
0.75 & 59.0 & 79.0 & \textbf{7.2} \\
\bottomrule
\end{tabular}
\caption{Layered TempoMem on Eventful, ViViT-B / Kinetics-400
(matching-frame accounting for all rows). Accuracy is flat for all
$\gamma \ge 0.25$ while compute falls $3.8\times$: TempoMem can discard up
to 75\% of Eventful's recompute set for a fixed 2-point Top-1 cost.
Whole-clip reference points under the same setting: dense 73.0\% Top-1 at
3355 GFLOPs/clip; canonical Eventful 71.0\% at 613 GFLOPs/clip (not
comparable to the matching-only rows above). Full-validation runs
(19{,}877 clips) confirm the endpoints: $\gamma = 0$ reaches 62.4\%
Top-1 / 82.2\% Top-5 and $\gamma = 0.5$ reaches 59.9\% / 79.8\% at the
same measured GFLOPs. \textbf{[TBD: full-validation runs for
$\gamma \in \{0.25, 0.75\}$ in progress.]}}
\label{tab:eventful-vivit}
\end{table}

\paragraph{Reading the accuracy-compute curve
(Fig.~\ref{fig:eventful-curve}).}
Figure~\ref{fig:eventful-curve} plots Top-1 against per-clip
matching-frame GFLOPs as the Eventful budget sweeps (green) and as
TempoMem's reuse fraction sweeps on top of the smallest Eventful budget
(blue). Two properties matter. First,
the TempoMem curve is a step, not a slope: the entire accuracy cost (2 points)
is paid at the smallest nonzero $\gamma$, after which accuracy is exactly
flat at 59.0/79.0 while compute falls from 20.6 to 7.2 GFLOPs/frame. This
says the tokens TempoMem drops beyond the first quartile were genuinely
redundant: Eventful flagged them as changed, but the change was already
represented in the prototype cache. Second, the step localizes the damage:
the loss comes from the first tranche of dropped tokens (those nearest the
matching threshold), suggesting an adaptive $\gamma$ could recover it.
Prototype granularity is irrelevant here: merge schedules of $T \in
\{2,4,8\}$ passes produced identical Top-1, Top-5, and FLOPs at every
$\gamma$, so on classification the filter is robust to how finely the cache
is compressed.

\begin{figure}[t]
\centering
%% PRELIM: figure legend uses "rho" for the reuse fraction; regenerate
%% with the paper's "gamma" symbol and full-validation numbers before
%% submission (the TempoMem name in the legend is already correct).
%% Source: results/comparison/plot_vivit_runs.log
\includegraphics[width=\linewidth]{fig_vivit_k400.pdf}
\caption{Top-1 accuracy vs.\ matching-frame GFLOPs per clip for ViViT-B /
Kinetics-400 (100 clips, Wilson 95\% CIs). Green: Eventful at increasing
token budgets. Blue: TempoMem's reuse-fraction sweep applied on top of the
smallest Eventful budget; the accuracy cost is paid entirely at the first
step and accuracy stays flat as compute continues to fall.}
\label{fig:eventful-curve}
\end{figure}

\subsection{Comparison with the State of the Art on ImageNet VID}
\label{sec:eval-vid}

\paragraph{Controlled comparison.}
The apples-to-apples evidence is Table~\ref{tab:vitdet-sweep}: Eventful and
the layered filter under the identical matching-frame protocol, differing
only in $\gamma$ (Eventful budget $k = 512$). Unlike classification,
detection accuracy falls monotonically with $\gamma$ (90.3\% to 59.3\%
mAP@50) as compute falls from 60.3 to 19.3 GFLOPs/frame; at $\gamma = 0.5$,
halving Eventful's compute costs 13 mAP points. Sweeping the Eventful
budget $k \in \{128, \dots, 1024\}$ shows the same picture along the whole
curve of this subset: at matched compute, Eventful reaches higher mAP@50
than the layered filter, and at matched mAP@50 Eventful is cheaper.
\textbf{On this paired sweep, the recompute-set filter does not improve
on Eventful's accuracy-compute frontier}; detection reuse is a genuine
trade of accuracy for compute rather than the near-free reuse observed on
classification. We report the trade-off as measured rather than tuning
around it, and analyze its cause below.

\paragraph{Full-dataset check.}
On all 639 VID videos, the layered filter with $\gamma = 0.5$ and
$k = 128$ reaches 60.5\% mAP@50 at 12.4 GFLOPs/frame under matching-frame
accounting, 5 points below the 25-video estimate for the same
configuration (65.5\%). \textbf{[TBD: the paired $\gamma = 0$ full run is
in progress; the full-dataset frontier comparison will be added when it
completes.]}

\begin{table}[t]
\centering
\small
\begin{tabular}{lccc}
\toprule
$\gamma$ & merge passes $T$ & mAP@50 (\%) & GFLOPs/frame \\
\midrule
0.00 (= Eventful) & -- & 90.3 & 60.3 \\
0.25 & 2 & 86.9 & 46.6 \\
0.25 & 4 & 86.4 & 46.6 \\
0.25 & 8 & 85.2 & 46.6 \\
0.50 & 2 & 77.8 & 32.9 \\
0.50 & 4 & 76.9 & 32.9 \\
0.50 & 8 & 74.7 & 32.9 \\
0.75 & 2 & 59.3 & 19.3 \\
\bottomrule
\end{tabular}
\caption{Layered TempoMem on Eventful, ViTDet-B / ImageNet VID
($k = 512$, matching-frame accounting for all rows). Detection accuracy
falls monotonically with the reuse fraction $\gamma$, and coarser
prototypes (more merge passes) hurt, the opposite of the classification
result. \textbf{[TBD: 25-video subset; evaluation on more videos in
progress.]}}
\label{tab:vitdet-sweep}
\end{table}

\paragraph{Reference frontier.}
Table~\ref{tab:vitdet-sota} places the compared methods under
whole-stream per-frame accounting (all frames, warm-up included),
together with their measured latency and peak memory. Eventful and STGT
define the current frontier; MaskVD follows; the KV-reuse-only standalone
variant sits well inside it. \textbf{[TBD: the layered TempoMem row
requires a whole-stream re-measurement; a frontier figure with the newest
runs will be added when it completes.]}

\begin{table}[t]
\centering
\resizebox{\columnwidth}{!}{%
\begin{tabular}{lcccc}
\toprule
Method & GFLOPs/fr. & Latency (ms) & Mem.\ (MB) & mAP@50 (\%) \\
\midrule
ViTDet-B (dense) & 174.5 & 69.6 & 760 & 84.1 \\
Eventful & 60.5 & 58.3 & 2118 & 84.0 \\
STGT & 68.5 & 56.2 & 1242 & 83.8 \\
MaskVD & 88.8 & 56.0 & 927 & 83.1 \\
TempoMem (standalone, KV-reuse only) & 156.5 & 101.8 & 929 & 72.0 \\
TempoMem (layered on Eventful) & \textbf{[TBD]} & \textbf{[TBD]} &
  \textbf{[TBD]} & \textbf{[TBD]} \\
\bottomrule
\end{tabular}}
\caption{ImageNet VID comparison with the state of the art under
whole-stream per-frame accounting and an identical GPU timing protocol.
All rows share the same accounting; the matching-only numbers of
Table~\ref{tab:vitdet-sweep} are not comparable to these and are therefore
kept separate.}
\label{tab:vitdet-sota}
\end{table}

% [TBD] A frontier figure (mAP@50 vs GFLOPs/frame, whole-stream
% accounting) will go here once the newest layered runs are
% measured whole-stream; the old fig_pareto_vid.png was removed because
% it predates them.

\subsection{When Does Content-Based Reuse Help?}
\label{sec:eval-analysis}

The ViViT and ViTDet results are opposite, and the contrast is the most
informative finding of this evaluation. On classification, the prediction
is read from a class token that aggregates the sequence; replacing a
redundant patch token with a nearby prototype changes what the class token
attends to only marginally, so accuracy is flat while 75\% of the
recompute set is dropped (Table~\ref{tab:eventful-vivit}). On detection,
the prediction is read from the spatial tokens themselves; every dropped
or substituted token is a location the detector can no longer describe
exactly, so mAP@50 falls in proportion to the reuse rate
(Table~\ref{tab:vitdet-sweep}). It is worth being precise about what
does and does not transfer between the two tasks: the \emph{compute}
side of the mechanism transfers unchanged, since the filter removes the
same fraction of the recompute set and the measured GFLOPs fall
proportionally on both benchmarks; what changes is the \emph{price} of
that reduction, from a small one-time accuracy step on action
recognition to a proportional mAP@50 loss on detection. The difference
is not the dataset but the readout. An action recognition model
answers one question per clip, and its class token aggregates evidence
over hundreds of tokens, so replacing a redundant background token with
a nearby prototype perturbs a sum in which that token is one term among
many. A detector answers a question at every location, and the feature
of each spatial token becomes the local evidence for a box at that
position, so the same perturbation lands directly on an output rather
than being averaged away. Two secondary observations support this
reading. First, prototype granularity is inert on classification but
matters on detection, where fewer merge passes (finer prototypes, hence
smaller substitution error per token) are consistently better. Second,
even the KV-reuse-only standalone variant, which keeps the token grid
fully intact and reuses nothing but the key and value projections, loses
12 mAP points on VID (Table~\ref{tab:vitdet-sota}): detection is
sensitive not only to dropped tokens but to approximated keys and values
themselves. TempoMem is therefore best positioned as an accelerator for
classification-style, global-readout video inference, and as a research
direction, rather than a finished method, for dense spatial prediction.

\subsection{Deployment Cost: Latency and Memory}
\label{sec:eval-deploy}

FLOPs do not automatically become speed, and we report this honestly.
Table~\ref{tab:vitdet-sota} includes measured wall-clock latency and peak
memory on ImageNet VID. The gated baselines translate their FLOPs savings
into real speedups (56 to 58\,ms vs.\ 69.6\,ms dense). TempoMem's
KV-reuse-only variant does not yet: matching and cache management add
overhead that outweighs the saved projections in the current
implementation (101.8\,ms), although
its memory footprint (929\,MB) stays close to dense and far below
Eventful's buffers (2118\,MB), reflecting the merged cache of
Proposition~\ref{prop:cachesize}. The gap is implementation-level, not
algorithmic: the matching product and gather/scatter steps are unfused
kernels, and profiling and fused implementations are future work.
\textbf{[TBD: latency and memory for the layered filter to be
measured under the same harness.]}

\subsection{Ablations}
\label{sec:eval-ablations}

\textbf{[TBD: ablations currently use the evaluation subsets of the
Experimental Setup; extended evaluation on more videos is in progress.]}

\paragraph{Reuse fraction $\gamma$.}
Covered by Tables~\ref{tab:eventful-vivit} and~\ref{tab:vitdet-sweep}:
a flat step on classification, a monotonic trade on detection.

\paragraph{Prototype granularity ($T$ merge passes).}
On ViViT, $T \in \{2,4,8\}$ is fully inert: identical accuracy and FLOPs
at every $\gamma$. On ViTDet, accuracy degrades with $T$ at every $\gamma$
(e.g., 77.8\% vs.\ 74.7\% mAP@50 at $\gamma = 0.5$), so detection prefers
the finest prototypes the memory budget allows. This is direct evidence
that classification tolerates coarse summaries while detection consumes
the per-token detail that merging removes.

\paragraph{Substitution on and off.}
The substitution variant of the layered mode refreshes each
matched candidate's K/V with its prototype instead of silently dropping
it. On classification it is accuracy-neutral: at $\gamma = 0.5$,
substitution on and off both score 59.0\% Top-1 and 79.0\% Top-5, while
substitution costs extra qkv compute (288 vs.\ 222 matching GFLOPs per
clip). On detection it buys a small gain (87.2\% vs.\ 86.9\% mAP@50 at
$\gamma = 0.25$), consistent with detection consuming the per-token
detail that a silent drop leaves stale. We therefore default to the
cheaper filter-only form on classification and note substitution as an
option where dense readouts matter.

\paragraph{Does content matching help without dropping tokens?}
We also tested a motion-compensated variant that matches stale tokens by
content but never drops them, hypothesizing robustness to large frame
strides. It is refuted: across strides $\{1, 4, 8, 16\}$ on ImageNet VID,
Eventful and the variant degrade identically (differences $\le 0.6$
mAP@50 points, within noise; e.g., 62.9\% vs.\ 62.3\% at stride 8). The
value of the prototype cache lies in \emph{removing} compute, not in
improving the reference features, which motivated the recompute-set filter
as the surviving design.

\paragraph{Warm-up length $W$ and refresh period $P$.}
\textbf{[TBD: sweeps over $W$ and $P$, including scene-cut stress tests,
are in progress; results will be reported here.]}

\subsection{Standalone TempoMem}
\label{sec:eval-standalone}

The standalone mode appears in this evaluation through its
KV-reuse-only form on ViTDet (Table~\ref{tab:vitdet-sota}).
\textbf{[TBD: full standalone classification results (saliency split with
token substitution) will be added here; the saliency split is being
retuned (current 100-clip runs reach 54.0\% Top-1 vs.\ 73.0\% dense) and
re-run at full validation scale.]}

\subsection{Summary of Findings}
\label{sec:eval-summary}

(1)~Layered on Eventful, TempoMem drops up to 75\% of the recompute
set on Kinetics-400 for a fixed 2-point Top-1 cost, cutting matching-frame
compute $3.8\times$ below Eventful itself, with the entire cost paid at
the first reuse step. (2)~On detection, the trade-off is much less
favorable: mAP@50 falls in proportion to the reuse rate because dense
spatial readouts consume exactly the per-token information that
prototype substitution removes, and on the paired 25-video sweep the
filter does not improve on Eventful's accuracy-compute frontier. This
scopes the method's near-free reuse to global-readout tasks. (3)~On the
full VID set, the filter at $\gamma = 0.5$, $k = 128$ scores 60.5\%
mAP@50 at 12.4 GFLOPs/frame, 5 points below its subset estimate.
(4)~FLOPs savings do not yet translate into latency, and memory stays
near dense. \textbf{[TBD: standalone classification results, whole-stream
measurements for the filter, the paired full-VID Eventful run, and the
$W$/$P$ sweeps are in progress.]}

\section{Discussion and Limitations}
\label{sec:discussion}

\paragraph{Content-based vs.\ position-based reuse.}
Temporal-difference methods, including the gates of every host the
layered mode can attach to, reuse computation per position and therefore
break under camera motion, where every position changes while scene
content largely does not. TempoMem matches by content, so a background
token that moved across the grid can still hit the cache. The two
signals are complementary, which is precisely what the layered mode
exploits: the host detects change, the memory recognizes recurrence.

\paragraph{Deployment on edge devices.}
The setting TempoMem is designed for is a camera that runs for hours on
hardware with a hard compute or energy ceiling, and two properties of
the design speak directly to it. First, the layered mode acts as a
\emph{compute dial}: when the deployment constraint is GFLOPs rather
than the last accuracy point, an operator can push the reuse fraction up
and trade a small, measured accuracy cost for a multiple reduction in
steady-state per-frame work, on top of whatever accelerator the platform
already ships. Second, the Persistent Scene Memory is the right shape
for long streams: its size is bounded independently of stream length
(Proposition~\ref{prop:cachesize}), its construction cost amortizes
toward zero as the stream grows (the Amortized stream cost paragraph),
and periodic refresh bounds staleness without ever storing more than the
prototype set. A static or slowly panning camera, the common case in
IoT installations, is exactly the regime where scene content recurs for
hundreds of frames and a persistent memory keeps paying out. The main
caveat is that our current implementation does not yet convert FLOPs
into wall-clock latency (Table~\ref{tab:vitdet-sota}): the matching and
gather/scatter steps are unfused kernels, and until they are fused the
measured savings are energy-proportional-compute savings rather than
speedups.

\paragraph{Saliency lag.}
The standalone foreground/background split consumes an attention map that
is one step stale, and at matching frames that map was itself computed on a
reduced sequence. We accept this lag because the split only needs to be
reliable for the background side, which is temporally stable by definition.
A foreground token mis-assigned to the background costs accuracy only if it
also matches the cache, while a background token mis-assigned to the
foreground is recomputed exactly and costs only efficiency; the fixed
foreground budget, recomputed every frame, bounds the damage of errors on
rapidly changing tokens.

\paragraph{Fixed reuse budgets.}
Top-count matching (and the fixed drop fraction $\gamma$ in the layered
mode) gives batch-uniform shapes and a predictable compute budget,
but forces the same reuse rate on easy and hard frames: a frame with
genuinely novel background is still forced to reuse a fraction $r$ of it,
with the refresh period $P$ acting as the safety net. Adaptive,
similarity-thresholded budgets are a natural extension we leave to future
work.

\paragraph{Failure modes.}
(i)~A scene cut between refreshes yields up to $P{-}1$ frames matched
against an obsolete cache; the entropy-based split mitigates this by
pushing novel content into the always-recomputed foreground, but
background-level errors persist until refresh. (ii)~Approximation error
from reusing merged, aging K/V accumulates with distance from the last
refresh, making $P$ the main accuracy-efficiency knob after $r$.
(iii)~Substitution is unsuitable for position-sensitive attention; the
KV-reuse-only variant covers this case at reduced savings. The
experiments quantify the third effect \textbf{[TBD: the $W$ and $P$
sweeps that quantify the first two are in progress]}.

\section{Conclusion}

We presented TempoMem, a training-free, memory-centric framework for
efficient video transformer inference. TempoMem clusters the temporally
stable tokens of a warm-up window into a bounded per-block Persistent
Scene Memory and, on later frames, reuses the cached key/value
projections of matched tokens while optionally substituting them with
their prototypes. The framework runs in two modes: a standalone
accelerator built around a foreground-preserving saliency split, and a
layered filter that attaches to any temporal token-reduction method,
strictly shrinking the set of tokens the host recomputes and reducing
exactly to the host when disabled. On classification, the evidence is
encouraging: layered on Eventful Transformers, TempoMem discards up
to 75\% of the recompute set on ViViT-B/Kinetics-400 for a fixed 2-point
Top-1 cost, cutting steady-state compute $3.8\times$ below Eventful
itself. On dense spatial prediction, the same mechanism trades accuracy
for compute much less favorably, with mAP@50 falling in proportion to the
reuse rate, and our analysis attributes this cleanly to the prediction
readout: global readouts absorb prototype substitution, dense readouts
consume exactly the per-token detail it removes. Wall-clock
latency also does not yet reflect the FLOPs savings, pending kernel
fusion. We see three directions as most promising: adaptive
similarity-thresholded reuse budgets in place of fixed fractions,
detection-aware reuse that protects tokens near object boundaries, and
fused matching/gather kernels to convert the measured FLOPs reductions
into latency reductions.

\bibliography{references}

\end{document}
